\documentclass[12pt]{article}
\usepackage[letterpaper,margin=1in]{geometry}
\usepackage[T1]{fontenc}
\usepackage[utf8]{inputenc}
\usepackage{lmodern,microtype,amsmath,amssymb,graphicx,booktabs,array,longtable,placeins,chngcntr}
\usepackage{setspace,float}
\graphicspath{{../04_Figures/}}
\usepackage[style=apa,natbib=true,backend=biber,sortcites=true,sorting=nyt]{biblatex}
\addbibresource{../06_Bibliography/aeri_refs_JAE_JBES.bib}
\usepackage[colorlinks=true,citecolor=blue,linkcolor=blue,urlcolor=black]{hyperref}
\usepackage{caption}
\captionsetup{font=small,labelfont=bf}
% JBES format: 12pt, one-inch margins, fully double-spaced body text
% (\setstretch{1.75} gives 26 lines of text on a full page).
\setstretch{1.75}
\renewcommand{\topfraction}{.9}\renewcommand{\bottomfraction}{.8}\renewcommand{\textfraction}{.07}\renewcommand{\floatpagefraction}{.85}\setcounter{topnumber}{3}\setcounter{totalnumber}{4}
\setlength{\emergencystretch}{2em}
\newcommand{\clip}{\operatorname{clip}}
\newcommand{\ind}{\mathbf{1}}
\newcommand{\exhibit}[4]{\begin{figure}[htbp]\centering\includegraphics[width=\textwidth,height=.40\textheight,keepaspectratio]{#1}\caption{#2}\label{#3}\begin{singlespace}\small\textit{Notes:} #4\end{singlespace}\end{figure}}
\begin{document}
\thispagestyle{plain}
\begin{center}
\begin{spacing}{1.15}
{\Large\bfseries Measuring Economic Risk in Africa: The African Economic Risk Index and African Risk Score\par}
\end{spacing}
\end{center}
\begin{center}\textbf{Abstract}\end{center}\vspace{-6pt}
\noindent We construct news-based measures of economic risk for 54 African economies from 2015 to 2025. The African Economic Risk Index (AERI) assigns graded scores to articles according to economic relevance, severity and textual emphasis and averages those scores over the monthly news flow. Eight non-exclusive pillar indices retain the domain information. The African Risk Score (ARS) maps bounded measures of AERI intensity, relevant-article breadth, average severity and active-domain count to a 0--100 scale. The data contain 21.30 million retained article records and 6,488 country-months that meet the principal descriptive coverage threshold. Exact-text deduplication removes 497,091 repeated records (2.33 percent) from the reconstructed frozen vintage; original and duplicate-adjusted AERI remain highly correlated (Pearson 0.9296; rank 0.9883). Aggregate AERI peaks in April 2020, but country histories differ in timing and composition. Alternative ARS weights produce closely aligned scores. In pooled panel VARs, inflation-risk readings precede higher headline inflation, exchange-rate depreciation is followed by higher AERI, and stronger real activity is followed by lower AERI. An external monetary-policy design also shows larger subsequent AERI increases after U.S. tightening shocks in economies with greater predetermined external-debt exposure. The reported sensitivity and falsification checks do not materially change these findings. AERI, the pillar indices and ARS are decomposable measures of reported economic risk that can be read alongside conventional macroeconomic statistics.

\begin{spacing}{1.15}
\noindent\textit{Keywords:} Alternative data; Newspapers; Panel VAR; Risk measurement; Text as data.\\
\textit{JEL Classification:} E37, C55, C81, O55.
\end{spacing}
\newpage


\section{Introduction}\label{sec:introduction}
Official statistics remain the main source of information about inflation, output, fiscal conditions and external balances, but they do not capture all changes in the economic environment as they occur. News reports often record shortages, financing pressures, currency stress, price increases and production disruptions before the full effects appear in conventional data. We develop a way to measure that information systematically across African economies.

We construct the \textbf{African Economic Risk Index (AERI)} from article-level scores that combine economic relevance, severity and textual emphasis. For each country and month, AERI averages these scores over all retained articles and multiplies the result by 100. Non-relevant articles remain in the denominator, so the index rises when economic-risk reporting becomes more prevalent, more severe or more prominent relative to the overall news flow.

We also construct eight non-exclusive pillar indices covering inflation, foreign-exchange and external conditions, debt and fiscal conditions, the financial sector, growth and the real economy, commodities and supply, political economy, and climate and disasters. These indices show where measured risk is concentrated. A separate PCA summary records common movement across the pillars. Finally, the \textbf{African Risk Score (ARS)} combines AERI intensity, breadth of relevant reporting, average severity and the number of active domains on a 0--100 scale. AERI is the principal measure; the pillars and ARS provide additional information about the composition and breadth of the underlying news signal.

The data contain 21.30 million retained article records for 54 African economies from January 2015 through December 2025. We observe 7,127 country-months, of which 6,488 contain at least 100 retained articles and form the principal descriptive sample. Aggregate AERI reaches its sample maximum in April 2020. The national series do not peak at the same time, and their risk composition differs across countries. ARS is also insensitive to an alternative weighting rule: the benchmark and PCA-weighted versions have a correlation of 0.995 in levels and 0.997 in ranks.

To examine economic content, we estimate pooled panel vector autoregressions (PVARs) linking the news measures to inflation, exchange-rate movements and monthly GDP growth. Inflation-risk lags foreshadow higher inflation, depreciation is followed by higher AERI, and stronger real activity is followed by lower AERI. We also use high-frequency U.S. monetary-policy surprises from \citet{jarocinski2020} as an external source of variation and interact them with 2014 external-debt exposure. AERI rises more after a tightening in African economies that entered the sample with greater external exposure. The differential response remains under alternative exposure measures, an alternative monetary-policy shock, leave-one-country-out estimation, and a sample restricted to FOMC announcement months.

The construction builds on the newspaper-based measurement approach of \citet{baker2016}, but the object is broader economic risk rather than policy uncertainty. Article scores are graded instead of binary. A single article may enter more than one economic domain. ARS is built separately from AERI using bounded inputs for intensity, breadth, severity and active domains. The same scoring rule is applied across 54 economies with very different media depth and economic structures, and coverage is reported explicitly.

The paper also relates to work in applied econometrics that treats news text as a source of macroeconomic information \citep{thorsrud2020,ellingsen2022,kalamara2022}. Those studies draw on deep archives from a single advanced economy. We take the same measurement problem to 54 economies whose archives differ in depth, language and continuity, and where the official series that text measures are usually checked against are often thin or late.

\section{Literature Review}\label{sec:literature}
\subsection{Economic risk, beliefs and news}
Economic decisions depend on prospective outcomes as well as current conditions. \citet{bloom2009} shows how uncertainty can delay investment and hiring, while \citet{jurado2015} measures uncertainty through the unpredictability of economic outcomes. A news-based risk index captures a different object: the salience and severity of economic concerns in reporting. A clearly anticipated increase in financing costs may generate substantial risk reporting even when there is little uncertainty about whether it will occur. Information also reaches households and firms unevenly \citep{mankiwreis2002}, and narratives shape how economic interpretations spread \citep{shiller2017}. News records changing attention and beliefs, but it also responds to realised conditions, so the empirical analysis examines movements from news to macroeconomic outcomes and feedback from realised conditions to the news measures. Section~S8 of the Supplementary Material gives the longer discussion.
\subsection{Newspaper-based indices of uncertainty and risk}
\citet{baker2016} construct economic policy uncertainty from articles jointly referring to the economy, policy and uncertainty. \citet{caldara2022} develop a geopolitical-risk measure, and \citet{ahir2022} use uncertainty-related language in country reports to extend international coverage. AERI retains the text-based approach but separates economic reporting into eight domains and incorporates severity and structural breadth. The domain structure shows whether reported stress is concentrated in particular parts of the economy rather than only whether a specified class of uncertainty terms appears.

For a text-based measure, the counted event and its denominator have to match the economic interpretation. AERI uses a graded article score, whereas ARS combines bounded information from several features of the monthly news flow. The eight pillars show where reported risk is concentrated. External-financing stress and domestic price pressure, for example, can rise at the same time without representing the same economic problem.

AERI also differs substantively from economic policy uncertainty. Its dictionary spans prices, production, finance and other economic domains, including developments whose occurrence may be relatively certain, and its article scores vary with severity and textual emphasis, whereas EPU is built around joint references to the economy, policy and uncertainty. The two share a news-based measurement tradition but represent different constructs and scales.
\subsection{News text as an econometric measurement problem}\label{sec:literature_text_econometrics}
A second body of work, much of it in applied econometrics and forecasting, treats newspaper text as data to be modelled and tested rather than only counted. \citet{thorsrud2020} extracts topics from a Norwegian business daily and combines them with quarterly GDP in a time-varying dynamic factor model. The resulting daily coincident index dates business-cycle phases closely and nowcasts output about as well as forecast combinations and expert judgement. \citet{larsen2019} use the same archive to show that news topics predict a range of macroeconomic and financial variables, and that an index built from the most informative topics is followed by persistent movements in financial markets, credit and borrowing. \citet{ellingsen2022} run a real-time forecasting comparison between 22.5 million Dow Jones articles and the FRED-MD database for US output, consumption and investment growth. The news data carry information that the hard indicators miss, most clearly for consumption.

The choice of text method has received as much attention as the choice of data. \citet{kalamara2022} compare dictionary counts with machine-learning methods applied to three UK newspapers. Text improves forecasts of GDP, inflation and unemployment, the gains are largest in periods of economic stress, and the best results come from term counts combined with supervised learning. \citet{barbaglia2023} score only the text that is grammatically linked to an economic term of interest, using a lexicon written for economics and finance, and find that the resulting sentiment improves forecasts of US macroeconomic variables, including the tails of the forecast distribution. \citet{shapiro2022} report that a lexicon designed for economic news outperforms general-purpose sentiment lists. \citet{bybee2024} estimate a topic model on the full text of the \textit{Wall Street Journal} and show that the attention paid to different themes closely tracks a wide range of economic activity. A related lesson runs through this work: word lists need to fit the domain. \citet{loughran2011} show that a general-purpose negative-word list classifies many words as negative that are not negative in a financial context.

These papers share two features. Each draws on deep archives from a single advanced economy, and the more flexible methods learn their topics or weights from that archive or from a target series observed with reasonable accuracy. Neither feature carries over to our sample. Archive depth, language and outlet continuity differ widely between countries, and for many of them the official series that a supervised model would be trained on arrive late, are revised heavily or are themselves interpolated. A topic model estimated country by country would also produce topics that differ from one country to the next, so a high reading in one country would not carry the same meaning as a high reading in another. AERI therefore applies one published dictionary and one scoring rule, with English, French and Portuguese terms, across all 54 economies. The graded severity tiers and textual-emphasis weights bring the score closer to the intensity measures used in the sentiment literature while keeping each article's contribution traceable to explicit rules.
\subsection{Composite indices, factor structure and missing observations}\label{sec:literature_composite}
Composite measures of economic conditions have a long history in applied econometrics. \citet{stockwatson2002} show that a small number of factors extracted from many series carry most of their forecasting content. \citet{mariano2003} and \citet{aruoba2009} build coincident indices from indicators observed at different frequencies, and \citet{lewis2022} construct a weekly index of US real activity from ten high-frequency series. \citet{banbura2014} extend maximum-likelihood factor estimation to data with an arbitrary pattern of missing observations, which is the usual situation when indicators start at different dates or are published with different delays.

Our aggregation step sits within this tradition but makes different choices for a different purpose. Principal components enter only after the article-scoring rules have fixed what counts as economic risk \citep{gentzkow2019,jolliffe2016}, and the weights come from covariance rather than from fit to inflation or output, so the composite is not fitted to the outcomes it is later compared with (Section~\ref{sec:ars}). Missing source months are kept as missing rather than filled in by a model. A factor model in the spirit of \citet{banbura2014} could estimate those months, but the published index would then depend on the model used to estimate them. Keeping the gaps visible leaves that choice to the user and keeps a missing archive month distinct from an observed month with no measured risk.


\subsection{African Economies and Alternative Sources of High-Frequency Information}

A growing literature uses alternative data where official statistics are sparse or late, including satellite imagery for local well-being across Africa \citep{yeh2020}, high-frequency search and payments data for Sub-Saharan activity \citep{buell2021}, and machine-learning nowcasts of GDP in data-constrained settings \citep{fotopoulou2026}. News data capture another part of the high-frequency measurement problem by recording economic concerns as they appear in public reporting. The text-based forecasting studies reviewed above use archives from the United States, the United Kingdom and Norway, where conventional data are already timely. The case for news-based measurement is stronger where official statistics are thinner, and that is also where the measurement problems are hardest. Cross-country comparability, which depends on publication volume, source composition, language and archive access, is handled through the scaling and coverage rules in Section~\ref{sec:methodology} and discussed further in Supplementary Section~S8. The pillar indices keep price pressures, fiscal concerns, external financing conditions and the other domains separate, including months when the aggregate index changes little. We compare the news measures with official macroeconomic statistics and examine whether the timing and sign of the relationships match the definitions of the measures.
\subsection{External Identification and Global Monetary Shocks}\label{sec:literature_external_id}
High-frequency monetary-policy identification lets us test AERI against a shock measured outside the news-scoring system. \citet{jarocinski2020} separate pure monetary-policy surprises from central-bank information shocks using the high-frequency co-movement of interest rates and equity prices around policy announcements. We interact their pure monetary-policy component with predetermined country exposure. The test asks whether the same global tightening produces a larger AERI response in countries that entered the sample with greater external financial exposure.

Inference in this design has to respect its structure. The shock varies over months and the exposure varies across countries, so errors may be correlated in both dimensions; two-way clustering by country and month allows for both \citep{cameron2011}. The number of country clusters is small, however, and \citet{mackinnon2017} show that conventional cluster-robust \textit{t}-tests can over-reject badly when clusters are few or differ widely in size, while the restricted wild cluster bootstrap generally performs better. We therefore read the exposure coefficient against finite-cluster reference distributions and report restricted wild-cluster-bootstrap results alongside them.
\section{Measuring AERI and ARS}\label{sec:methodology}
We build the measures from the collected articles, score each article, and then aggregate the scores by country and period. The output consists of AERI, eight overlapping pillar indices, an eight-pillar PCA common-factor composite, and the bounded ARS. Table~\ref{tab:measure_map} gives the definition and role of each measure.
\begin{table}[htbp]\centering\small\singlespacing
\caption{Measurement objects and their roles}\label{tab:measure_map}
\setlength{\tabcolsep}{4pt}
\begin{tabular}{p{2.5cm}p{4.7cm}p{6.1cm}}\toprule
Measure & Construction & Interpretation\\\midrule
AERI & $100$ times the sum of article-level risk scores divided by all retained articles in a country-month & Overall intensity of scored economic-risk content relative to the observed news flow; the principal economic-risk index\\
Pillar index & Same article-level risk score as AERI, with the numerator restricted to articles assigned to one of eight economic domains & Domain-specific diagnostic showing where measured economic risk is occurring; overlapping assignments preserve cross-domain connections\\
Eight-pillar PCA composite & Log-transformed, percentile-scaled pillar indices combined with fixed first-component weights & Common-factor diagnostic showing how strongly the eight risk domains move together; useful for distinguishing broad cross-domain stress from a concentrated pillar episode\\
ARS & Weighted average of bounded AERI intensity, relevant-article breadth, conditional severity and structural spread, scaled to 0--100 & Multidimensional monitoring score summarising how intense, widespread, severe and broad-based the measured risk is\\
\bottomrule\end{tabular}
\par\medskip\raggedright\footnotesize\textit{Notes:} AERI is the principal article-score index. The pillar indices locate the economic domains in which that risk appears. The PCA composite provides a companion diagnostic of common cross-domain movement, while ARS is a separate bounded monitoring score. Each measure is interpreted according to its stated construction and role in the framework.
\end{table}

\subsection{Data Collection}\label{sec:data_collection}
We assemble monthly country-level records from national news sources and organise retrieval by country, site and month, with bounded retrieval attempts and site-level diagnostics. The recorded retrieval settings allow the collection process to be repeated, although outlet coverage can still vary and repeated stories may remain. Corpus size is reported in retained article records. Section~\ref{sec:deduplication} evaluates the sensitivity of AERI and ARS to repeated identical article bodies using a separately reconstructed, internally matched archive.

Our analytical vintage contains 21,303,573 article records across 54 African countries from January 2015 to December 2025. Of these, 5,896,776 meet the economic relevance criterion. The 7,127 country-month records include 639 months with fewer than 100 articles. To limit sensitivity to very small news denominators, the principal descriptive comparisons use the remaining 6,488 country-months with at least 100 retained articles.\footnote{Months below the threshold remain in the underlying series and coverage files. The threshold governs the principal cross-country descriptive comparisons and the ARS calibration sample, while preserving the lower-volume observations in the documented archive.} The median country-month contains 1,039 articles. Months for which the source archive is absent or unreadable are recorded as missing. Supplementary Section~S11 describes the collection architecture.

The country-month is the unit of aggregation. A widely syndicated event can therefore generate several article records, sometimes across multiple outlets. We keep record counts distinct from independent stories and apply the coverage threshold at the country-month level before making cross-country descriptive comparisons. Calibration, descriptive analysis and macroeconomic estimation use separately documented samples because their data requirements differ. Section~\ref{sec:deduplication} quantifies how much the published series changes when identical article bodies are collapsed within country-months.
\exhibit{Figure_1.pdf}{News coverage and eligible country observations}{fig:coverage}{The left panel reports total retained article records by month. The right panel counts countries with at least 100 articles. Article counts refer to retained records before story-level exact-text deduplication. Changes in the archive can affect measured coverage.}
\subsection{Data Pre-processing}\label{sec:preprocessing}
Preprocessing converts article text into economic relevance, severity, textual emphasis and pillar assignments. The same rules apply across the retained corpus. An article can activate several pillars because economic developments often span more than one domain.
\subsubsection{Text Normalisation}\label{sec:text_norm}
The headline and body form a single textual unit, $x_i$. Normalisation removes formatting differences before dictionary matching. Scoring is rule-based and uses observable textual features from this normalised text.
\subsubsection{Economic Relevance Filtering}\label{sec:relevance_filter}
Relevance is defined by
\begin{equation}R_i=\ind\{x_i\text{ matches the economic-relevance dictionary}\}.\label{eq:relevance}\end{equation}
Non-relevant articles receive zero economic-risk score but remain in the denominator. Relevant reporting also includes discussions of economic conditions expressed without acute-crisis language.

All retained articles enter the denominator, so the share of the news flow devoted to economic risk remains part of the signal. The separate breadth and severity diagnostics then recover, respectively, the prevalence of relevant reporting and its conditional intensity. For an observed source month with retained records, a zero score is therefore an empirical outcome; missing source months are handled separately.
\subsubsection{Severity Tiering}\label{sec:severity_tiering}
Let $\mathcal K_i\subseteq\{1,2,3,4\}$ contain the severity tiers matched in the text. We set
\begin{equation} S_i=\begin{cases}0,&R_i=0,\\ \max(\{1\}\cup\mathcal K_i),&R_i=1.\end{cases}\label{eq:severity}\end{equation}
A relevant article receives a baseline of one, with stronger matched language increasing the score. Where several cues occur, the highest tier applies. The resulting severity tier is an ordinal description of language strength.
\subsubsection{Textual Emphasis Weighting}\label{sec:confidence_weighting}
The textual weight combines a headline cue $T_i$, a body-length indicator $L_i$ for at least 1,200 normalised characters, and an indicator $M_i$ for at least two distinct active pillars:
\begin{equation}W_i=\clip(1+0.5T_i+0.2L_i+0.3M_i,1,2).\label{eq:confidence}\end{equation}
Here $\clip(x,a,b)=\min(\max(x,a),b)$. The weight is a bounded textual-emphasis adjustment based on placement, length and cross-domain content. Its coefficients belong to the article-scoring rule and are distinct from the PCA aggregation weights used later in the construction.
\subsubsection{Article-Level Risk Scoring}\label{sec:article_scoring}
The score is
\begin{equation} A_i=R_iS_iW_i,\qquad 0\leq A_i\leq8.\label{eq:article}\end{equation}
Each activated pillar contributes once to an article score; repeated phrases within a story therefore leave that pillar's contribution unchanged. Supplementary Table~S9 shows representative terms from the production dictionaries and a worked example of how the matching rules map one article into AERI, pillar and ARS inputs.
\subsection{Construction of the African Economic Risk Index}\label{sec:aeri}
For country $c$ and month $t$, let $\mathcal A_{ct}$ denote the retained articles and $N_{ct}=|\mathcal A_{ct}|$. AERI is
\begin{equation}\mathrm{AERI}_{ct}=\frac{100}{N_{ct}}\sum_{i\in\mathcal A_{ct}}A_i.\label{eq:aeri}\end{equation}
The index measures scored economic content per article. It can exceed 100, with a theoretical upper bound of 800. Annual AERI weights monthly values by article counts, so the annual numerator and denominator correspond to the same article universe.

Define $\overline{SW}_{ct}^{\,R}=\sum_iR_iS_iW_i/\sum_iR_i$ whenever relevant articles are present. AERI can then be written as
\begin{equation}
 \mathrm{AERI}_{ct}=100B_{ct}\overline{SW}_{ct}^{\,R}.\label{eq:aeri_decomposition}
\end{equation}
If there are no relevant articles, AERI is zero for a non-empty month. Equation~\eqref{eq:aeri_decomposition} shows why rising AERI can reflect a larger share of economic-risk reporting, more severe or emphatic relevant reporting, or both. The decomposition also shows that breadth and intensity are related inputs to ARS. PCA summarises their covariance rather than treating the dimensions as statistically independent.
\subsection{Pillar-Level Decomposition and Diagnostics}\label{sec:pillars_diagnostics}
Let $I_{pi}$ indicate that article $i$ activates pillar $p$. The pillar index is
\begin{equation}P_{pct}=\frac{100}{N_{ct}}\sum_{i\in\mathcal A_{ct}}A_iI_{pi}.\label{eq:pillar}\end{equation}
The eight pillars can overlap.\footnote{This treatment parallels the category-specific newspaper indexes in \citet{baker2016}, where a single article may enter more than one category.} Each pillar lies between zero and AERI, and the sum can exceed AERI when a story contributes to several domains. Breadth, severity and structural coverage are
\begin{equation}B_{ct}=\frac{\sum_iR_i}{N_{ct}},\quad V_{ct}=\frac{\sum_iS_i}{\max(1,\sum_iR_i)},\quad K_{ct}=\sum_{p=1}^{8}\ind\{P_{pct}>0\}.\label{eq:diagnostics}\end{equation}
Each raw pillar index measures one domain. To summarise their common cross-domain movement, define $\ell_{pct}=\log(1+P_{pct})$ and
\begin{equation}Q_{pct}=\clip\left(\frac{\ell_{pct}-a_p}{b_p-a_p},0,1\right),\quad C_{pct}=100u_pQ_{pct},\quad H_{ct}=\sum_{p=1}^{8}C_{pct}.\label{eq:pca_pillars}\end{equation}
The bounds $a_p$ and $b_p$ are the first and ninety-ninth calibration percentiles, and $u_p$ are fixed weights obtained from the first principal component of the standardised transformed pillars. The resulting $H_{ct}$ is the eight-pillar PCA composite, which we use as a common-factor diagnostic. AERI can be high when overall scored risk content is elevated but concentrated in a few domains; the PCA composite rises most strongly when several transformed pillar indices move together. The comparison helps distinguish broad cross-domain stress from episodes driven mainly by one or two pillars.

The PCA composite captures cross-domain co-movement that $K_{ct}/8$ does not. The ARS component $K_{ct}/8$ counts active pillars; $H_{ct}$ also incorporates their transformed intensity and empirical common-factor weights. Supplementary Table~S10 reports the corresponding loadings, scaling bounds and weights.

\subsection{Construction of the African Risk Score}\label{sec:ars}
ARS combines four bounded inputs:
\begin{equation}X_{ct}=\left(\clip\left(\frac{\mathrm{AERI}_{ct}}{200},0,1\right),B_{ct},\frac{V_{ct}}4,\frac{K_{ct}}8\right)',\qquad\mathrm{ARS}_{ct}=100\sum_{j=1}^{4}w_jX_{jct}.\label{eq:ars}\end{equation}
The divisor 200 sets the saturation point for the intensity component. The four entries have direct interpretations: $\clip(\mathrm{AERI}/200,0,1)$ is bounded AERI intensity, $B$ is the share of retained articles that are economically relevant, $V/4$ scales conditional severity to the unit interval, and $K/8$ records the share of the eight domains that are active. ARS is therefore a bounded 0--100 summary of the reported risk environment, rising as risk-related reporting becomes more intense, prevalent, severe or structurally widespread.

The four components are allowed to covary. Breadth is already one margin of AERI through equation~\eqref{eq:aeri_decomposition}, but it enters ARS separately by design. This allows two country-months with similar AERI to receive different ARS values when one reflects narrow, highly intense reporting and the other reflects broader, multi-domain stress. PCA summarises the covariance among these diagnostics and reduces the dependence of the final score on arbitrary hand-set weights.

Standardise the calibration inputs as $Z_j=(X_j-\mu_j)/s_j$, using population standard deviations. Let $v$ be the unit-length first principal-component loading vector, oriented toward higher component values. The weights on the bounded input scale are
\begin{equation}v=\arg\max_{\|b\|=1}b'\widehat\Sigma_Zb,\qquad w_j=\frac{v_j/s_j}{\sum_kv_k/s_k}.\label{eq:weights}\end{equation}
Dividing by $s_j$ converts standardised loadings to coefficients on the bounded input scale. With positive loadings, ARS is a rescaled first component and is weakly increasing in its bounded inputs. The same conversion is applied separately to the transformed pillar vector $Q$ to obtain the fixed common-factor weights $u_p$ used in equation~\eqref{eq:pca_pillars}. Calibration rejects near-constant inputs or materially mixed-sign loadings rather than converting negative loadings to absolute values.

We calibrate ARS using 6,238 country-months with at least 100 articles during January 2015--July 2025. ARS weights are 26.25 percent for intensity, 27.22 percent for breadth, 32.23 percent for severity and 14.30 percent for structure. PC1 explains 62.22 percent of standardised variation in the four ARS inputs and 52.50 percent for the transformed eight-pillar system. Weights remain fixed between recalibrations. Historical scores remain retrospective because calibration uses information through July 2025. In the eligible 2015--2025 sample, the median month activates seven of the eight pillars, while the median relevant-article share is 25.8 percent. These diagnostics help explain why ARS can move differently from AERI even though both are generated from the same underlying article corpus.

Table~\ref{tab:component_weights_full} and Supplementary Table~S10 report the two PCA calibrations. Severity receives the largest ARS weight even though intensity has the largest standardised loading because equation~\eqref{eq:weights} converts standardised loadings back to coefficients on the bounded input scale before normalisation. Supplementary Table~S10 documents how the transformed domain indices load on their common factor and makes the construction used in Supplementary Figure~S5 auditable. Debt and fiscal risk and financial-sector risk receive the largest common-factor weights, at about 17.3 percent each, while commodity and supply risk receives 6.48 percent. Weights are held fixed between recalibrations, and a recalibration is treated as a new measurement vintage (Supplementary Section~S9).

\begin{table}[htbp]\centering\small\singlespacing
\caption{ARS component weights and standardised loadings}\label{tab:component_weights_full}
\setlength{\tabcolsep}{4pt}
\begin{tabular}{lrrr}\toprule
Component & Benchmark (\%) & PCA (\%) & PC1 loading\\\midrule
Intensity & 40.00 & 26.25 & 0.5803\\
Breadth & 25.00 & 27.22 & 0.5502\\
Severity & 20.00 & 32.23 & 0.3921\\
Structure & 15.00 & 14.30 & 0.4548\\
\bottomrule\end{tabular}
\par\medskip\raggedright\footnotesize\textit{Notes:} The 6,238 calibration observations cover January 2015--July 2025 and have at least 100 articles. PC1 explains 62.22 percent of standardised variation. Weights on the bounded input scale differ from standardised loadings; equation~\eqref{eq:weights} gives the conversion.
\end{table}
\section{Empirical Properties and Evaluation}\label{sec:results}
\subsection{Aggregate Risk Dynamics}\label{sec:aggregate_risk}
Figure~\ref{fig:trends} presents the two principal measures used throughout the paper. The left panel shows AERI as defined in equation~\eqref{eq:aeri}; the right panel shows ARS from equation~\eqref{eq:ars}, using the fixed PCA weights in Table~\ref{tab:component_weights_full}. For each month, the figure plots the unweighted cross-country median with interquartile ranges as shaded bands. The eight-pillar PCA composite is analysed separately in Supplementary Figure~S5.

The monthly median AERI reaches its sample maximum of 68.94 in April 2020, coinciding with the acute phase of the COVID-19 pandemic and the associated economic disruption. This timing provides a simple check on the economic content of the index. Country peaks occur in different months, however, and reading the aggregate series with country paths and pillar composition shows whether a spike reflects a broad continental episode or a more localised event. The shaded bands show dispersion across eligible countries. Because the set of eligible countries changes with archive coverage, aggregate movements are read alongside Figure~\ref{fig:coverage}.

Across 6,488 eligible country-months, AERI has a mean of 47.89 and a standard deviation of 27.56. ARS has a mean of 36.59 and a standard deviation of 9.66. Their different scales follow directly from construction: AERI preserves article-score intensity, while ARS combines bounded components. The two measures are read jointly, with AERI describing the level of the scored signal and ARS its broader configuration.
\exhibit{Figure_2.pdf}{Aggregate AERI and ARS dynamics}{fig:trends}{Lines are unweighted country medians; shaded regions are interquartile ranges. The annotation marks the COVID-19 pandemic shock at the maximum monthly median AERI in the sample (April 2020, 68.94). Only country-months with at least 100 articles enter. Neither series is weighted by GDP. ARS uses the fixed PCA calibration.}



Supplementary Figure~S5 compares AERI with the eight-pillar PCA composite and the benchmark with the PCA-weighted ARS. The AERI and composite medians correlate 0.558, the two ARS medians correlate 0.975, and all four series peak in April 2020. Outside that episode the composite separates broad cross-domain stress from concentrated spikes; May 2022, for example, is its second-highest month but not among the five highest AERI months.

\begin{table}[htbp]\centering\small\singlespacing
\caption{Distribution of AERI and alternative ARS specifications}\label{tab:comparison}
\setlength{\tabcolsep}{4pt}
\begin{tabular}{lrrrrr}\toprule
Measure & Mean & SD & Median & Minimum & Maximum\\\midrule
AERI & 47.89 & 27.56 & 46.91 & 0.00 & 300.16\\
ARS (PCA weights) & 36.59 & 9.66 & 37.27 & 0.00 & 85.37\\
Benchmark ARS & 35.42 & 10.71 & 36.01 & 0.00 & 88.93\\
\bottomrule\end{tabular}
\par\medskip\raggedright\footnotesize\textit{Notes:} Common sample: 6,488 country-months with at least 100 articles, 2015--2025. The benchmark and PCA-weighted ARS have Pearson correlation 0.9950, rank correlation 0.9971, and mean absolute difference 1.49 points. The comparison changes only the ARS weights; AERI itself has no PCA weighting stage.
\end{table}

Table~\ref{tab:comparison} shows the same pattern as Supplementary Figure~S5. The PCA-weighted and benchmark ARS have very similar central tendencies and nearly identical rankings, although the PCA specification is modestly less dispersed (standard deviation 9.66 versus 10.71). Changing the weights shifts the scale slightly while leaving the relative ordering of country-months largely intact. AERI is considerably more dispersed because it preserves the original article-score intensity and remains on its own scale.
\subsection{Annual Africa Economic Risk Indices by Country (2015--2025)}\label{sec:aeri_heatmap}
Supplementary Figure~S6 follows annual AERI for ten illustrative economies and shows that the continental median hides marked country differences in levels, persistence and the timing of peaks.
Among countries with at least nine eligible months in 2025, the highest annual AERI values are Zambia (112.40), Liberia (99.59), Lesotho (91.14). Annual values are article-count-weighted averages of monthly AERI scores. The resulting ordering describes the intensity of observed economic-risk reporting across the retained news archive.
We complement the annual comparison with exact monthly summary statistics for the same ten-country group in Supplementary Table~S8; the full all-country summary and monthly country paths are reported in the accompanying Supplementary Material. These outputs document differences in levels, variability and the timing of peaks that the aggregate median and annual averages conceal. Within the ten-country group, Nigeria has the highest mean monthly AERI, while Ethiopia exhibits the greatest time-series dispersion. These country results are descriptive and use the same AERI definition as the aggregate analysis.


\subsection{Pillar-Level Decomposition}\label{sec:pillar_results}
The pillar indices decompose the aggregate signal by economic domain. Supplementary Figure~S7 reports the pooled mean of each raw pillar index across the 6,488 eligible country-months, and Supplementary Table~S11 shows how the domain profile changes across broad subperiods. Because an article may activate several domains, these are overlapping measures of reported risk content rather than shares of a fixed total.

Financial-sector risk has the largest pooled mean, followed by debt and fiscal risk. FX/external, political-economy and inflation risks form a middle group; growth and real-economy risk sits below them, while climate/disaster and commodity/supply risks are the least prominent on average. Dictionary size does not account for the financial-sector ranking: that pillar contains nine top-level terms or phrases, the same number as six other pillars, while commodity/supply contains seven. Term frequency, outlet specialisation and the newsworthiness of particular events can differ across domains. For that reason, the level comparisons are descriptive, and the within-pillar changes over time carry more weight in interpretation. Those changes are informative: growth and real-economy risk rises from 3.14 in 2015--2019 to 4.85 in 2020, and inflation risk rises from 3.47 before the pandemic to 5.33 in 2021--2025 (Supplementary Table~S11).


\subsection{From AERI to ARS: Intensity, Breadth, and Structure}\label{sec:aeri_vs_ars}
ARS separates reporting intensity from the other components of the monthly news signal. Structural breadth is the number of domains with positive scores and can remain high even when AERI is low. Comparing the two measures shows whether an AERI increase is concentrated in a small number of domains or accompanies wider economic coverage.

Table~\ref{tab:comparison} shows that the broad ARS pattern is highly insensitive to the weighting choice: the 40/25/20/15 benchmark has a pooled correlation of 0.995 with PCA-weighted ARS. The pooled correlation between AERI and PCA-weighted ARS is lower, at 0.911. They clearly share a common signal, but ARS changes the scale and gives explicit weight to breadth, severity and structural spread.

\subsection{Evaluation of AERI and ARS}\label{sec:validation}
Following \citet{baker2016}, evaluation starts with the accounting identities and the coverage record. We then examine duplicate stories and alternative ARS weights, compare the series with salient economic episodes, estimate selected PVAR relationships, and use externally identified U.S. monetary-policy shocks interacted with predetermined country exposure. Each exercise addresses a different part of the measurement problem.

\subsubsection{Accounting, coverage and aggregation robustness}
We reconcile AERI directly to its article-score numerator and total-article denominator. Equation~\eqref{eq:aeri_decomposition} separates the relevance-share margin from conditional severity and emphasis. Pillar scores use the same article score with overlapping domain indicators, and ARS combines four bounded diagnostics under a frozen calibration. The published country-month measures can be traced back to those article-level calculations.

Coverage also matters. The 2015--2025 vintage contains 7,127 country-month records, of which 6,488 meet the 100-article threshold used for the principal descriptive comparisons. Lower-volume months remain in the documented archive but are omitted from those comparisons because ratio-based measures become unstable with very small denominators. Figure~\ref{fig:coverage} shows the resulting variation in the observation base. As an archive-depth diagnostic, we regress AERI on log retained article volume within the eligible sample. The within-country correlation is 0.035; the article-volume coefficient is 1.14 with country fixed effects ($p=0.702$) and $-1.47$ after adding year-month fixed effects ($p=0.640$). Supplementary Table~S12 reports the specifications. The estimates show little association between AERI and archive size within countries. Changes in source composition and reporting practice remain separate concerns.

We assess ARS weight sensitivity by replacing the PCA weights with the benchmark 40/25/20/15 rule. The broad score pattern is almost unchanged: Pearson correlation is 0.9950, rank correlation is 0.9971 and the mean absolute difference is 1.49 points. With the sample and component definitions held fixed, the exercise attributes the remaining differences to the weighting choice.

\subsubsection{Syndicated stories and exact-text deduplication}\label{sec:deduplication}
Repeated publication can amplify a country-month signal when the same high-scoring text appears in several outlets. We collapse identical decoded article bodies within each country-month, retaining one representative per duplicate group.\footnote{Exact matching is applied to decoded bodies of at least 200 characters. Appearances in different countries or months remain separate news exposures, while minor edits, translations and rewritten versions remain distinct records.}

The correction is applied algebraically to the frozen analytical vintage. Each removed record is scored with the original dictionaries, and its contribution is subtracted from the archived numerator, relevance and severity counts and each pillar numerator. The all-article denominator is reduced by the same record count. Source files are read-only, and duplicate counts, representative examples and file fingerprints are retained in the replication extension.

Across 21,303,020 frozen-vintage records represented in 7,126 reconstructed country-months, the rule removes 497,091 repetitions (2.33\%), including 186,379 cross-outlet copies. The reconstruction is 553 records and one Togo country-month short of the archived aggregate because 42 supplied source files were empty or unreadable. For months with a partial count mismatch, unmatched frozen-vintage records remain in the denominator and only directly observed duplicates are removed. On the common coverage sample, original and duplicate-adjusted AERI have Pearson correlation 0.9296 and rank correlation 0.9883. The mean absolute change is 2.32 AERI points, and the 95th-percentile absolute change is 8.30. Fixed-weight ARS has correlation 0.9680 and mean absolute change 0.61 points. 69 country-months that originally met the 100-article rule fall below it after removal. Supplementary Table~S2 reports the full comparison.


Country-level results show where duplication matters most. The median within-country AERI correlation is 0.9992, 48 of 53 countries retain the same peak month on the common sample, and the cross-country median still peaks in April 2020. A small number of country-specific spikes are sensitive to identical reprints, most clearly Sierra Leone in November 2016 (Supplementary Section~S4). The PVAR and external-shock specifications continue to use the original benchmark vintage.

\subsubsection{PVAR-based external-consistency evidence}
We use pooled PVARs and impulse-response functions to compare the news measures with subsequent macroeconomic outcomes. Inflation-risk reporting is paired with realised inflation, exchange-rate innovations with AERI, and monthly GDP growth with AERI. Table~\ref{tab:selected_external_consistency} reports these relationships; sample selection, lag choice, stability diagnostics and inference are documented in Supplementary Section~S3. All three joint lag tests are statistically significant at the paper's 10 percent criterion, with the inflation-risk and depreciation relationships also significant at 5 percent. At the reported horizons, the pointwise 90 percent intervals exclude zero in the predicted direction.\footnote{The bootstrap intervals are pointwise across horizons. They quantify horizon-specific uncertainty; simultaneous coverage of an entire impulse-response path would require a different band construction.}
\begin{table}[htbp]\centering\small\singlespacing
\caption{Directional PVAR external-consistency tests}\label{tab:selected_external_consistency}
\setlength{\tabcolsep}{4pt}
\begin{tabular}{@{}p{4.25cm}crp{6.0cm}@{}}\toprule
Relationship & Expected & Lag-test $p$ & Key response (90\% interval)\\\midrule
Inflation risk $\rightarrow$ headline inflation & $+$ & 0.0299 & M6: $+0.531$ pp $[0.287,\ 0.825]$\\
Depreciation $\rightarrow$ AERI & $+$ & 0.0142 & M2: $+0.0166$ SD $[0.0060,\ 0.0241]$\\
GDP growth $\rightarrow$ AERI & $-$ & 0.0624 & M1: $-0.181$ SD $[-0.338,\ -0.068]$\\
\bottomrule\end{tabular}
\par\medskip\raggedright\footnotesize\textit{Notes:} Expected signs are implied by the economic interpretation of the measures rather than chosen from the estimated responses. All three joint lag tests are statistically significant at the paper's 10 percent criterion, and the reported pointwise 90 percent intervals exclude zero in the expected direction. The inflation-risk and depreciation lag tests are also significant at 5 percent. Exact $p$-values are shown so readers can apply alternative significance conventions. These are conditional dynamic associations under the stated recursive ordering; structural identification is examined separately with the external monetary-policy shock design.
\end{table}


The strongest forward-looking result concerns the inflation-risk pillar. In the GDP-augmented PVAR, the three inflation-risk lags are jointly significant in the headline-inflation equation ($p=0.0299$). A one-within-country-standard-deviation inflation-risk innovation raises headline inflation by 0.111 percentage points on impact, 0.288 after one month and 0.392 after two months. The response continues to build and reaches 0.531 percentage points in month six, with a 90 percent interval of $[0.287,0.825]$. The pointwise 90 percent band stays above zero at every horizon from impact through month twelve. The hump-shaped profile indicates that the inflation-risk pillar contains information about sustained price pressure beyond the contemporaneous inflation reading. In the terminology of \citet{baker2016}, the news-based innovation \emph{foreshadows} subsequent inflation.
\exhibit{Figure_6.pdf}{Headline inflation response to an inflation-risk innovation}{fig:inflation_irfs}{The figure uses the 36-country GDP-augmented PVAR(3), estimated on 4,286 observations. The innovation is a one-within-country-standard-deviation increase in the article-score inflation-risk pillar. Bands are pointwise 90 percent intervals from 400 country-cluster bootstrap draws. The three inflation-risk lags are jointly significant in the headline-inflation equation ($p=0.0299$); the response is positive at every reported horizon and peaks at 0.531 percentage points in month 6.}

Figure~\ref{fig:aeri_macro_feedback} examines feedback from realised macroeconomic conditions to AERI. The exchange-rate result is particularly clear. A one-percentage-point innovation in monthly depreciation raises AERI by 0.0098 within-country standard deviations on impact, with a 90 percent interval of $[0.0033,0.0153]$. After a small month-one response, AERI rises to 0.0166 standard deviations in month two, with a 90 percent interval of $[0.0060,0.0241]$. The three depreciation lags are jointly significant in the AERI equation ($p=0.0142$). The response is short-lived rather than permanent. Its sign and timing match the interpretation of AERI as a measure that can register currency stress as exchange-rate pressure enters the news.

The real-activity response provides a complementary sign test. A one-percentage-point innovation in monthly GDP growth has a small and imprecise impact response, but AERI falls sharply one month later by 0.181 within-country standard deviations, with a 90 percent interval of $[-0.338,-0.068]$. Point estimates remain negative thereafter and gradually return toward zero. The joint three-lag Wald test is $p=0.0624$, which is statistically significant at the 10 percent level. The horizon-specific response and the joint lag test therefore point in the same direction: stronger activity is followed by less intense economic-risk reporting once information about improved real conditions is absorbed into the news flow.

Across the three IRFs, inflation-risk reporting precedes higher realised inflation, currency depreciation is followed by higher AERI, and stronger real activity is followed by lower AERI. The forward inflation system places the risk measure first; the feedback system places the macroeconomic variables first. The same sign pattern appears under both recursive orderings. We interpret the estimates as conditional dynamic relationships. Structural identification is handled separately in the external-shock exercise below.
\exhibit{Figure_7.pdf}{AERI responses to exchange-rate and real-activity innovations}{fig:aeri_macro_feedback}{Reverse 36-country GDP-augmented PVAR(3), estimated on 4,286 observations. Depreciation raises AERI, peaking at 0.0166 within-country standard deviations in month 2; GDP growth is followed by a $-0.181$ standard-deviation AERI response in month 1. Shading denotes pointwise 90 percent country-cluster bootstrap intervals. The opposite signs are consistent with higher reported risk under currency stress and lower reported risk after stronger activity.}

\FloatBarrier
\subsubsection{Externally identified monetary-policy exposure test}\label{sec:external_identification}
To complement the recursive PVAR, we turn to a shock measured outside the AERI system. Following \citet{jarocinski2020}, we use the pure U.S. monetary-policy surprise extracted from high-frequency movements in interest rates and stock prices around Federal Reserve announcements. We aggregate the surprise to the month and standardise it over 2015--2025. Country exposure is fixed at 2014 external debt stocks as a share of GNI from the World Bank International Debt Statistics \citep{worldbankdebt}, so the exposure measure predates the AERI sample.

For country $i$ and horizon $h$, the preferred local projection is
\begin{equation}
\begin{aligned}
\mathrm{AERI}_{i,t+h}-\mathrm{AERI}_{i,t-1}
={}&\alpha_i^h+\gamma_t^h+\beta_h\left(MP_t\times E_i^{2014}\right)\\
&+\sum_{\ell=1}^{3}\theta_{h\ell}\Delta\mathrm{AERI}_{i,t-\ell}
+\sum_{\ell=1}^{3}\psi_{h\ell}\left(MP_{t-\ell}\times E_i^{2014}\right)
+\varepsilon_{i,t+h}.
\end{aligned}
\label{eq:external_id}
\end{equation}
where $MP_t$ is a one-standard-deviation pure monetary-policy tightening and $E_i^{2014}$ is standardised $\log(1+\text{external debt/GNI})$. Country fixed effects absorb persistent differences in average risk reporting, while year-month fixed effects absorb the common global shock, COVID-19 and other disturbances shared across countries in a given month. The identifying variation is therefore cross-country exposure to the same externally measured tightening. Standard errors are two-way clustered by country and month following \citet{cameron2011}. We use finite-cluster reference inference and report restricted wild-cluster-bootstrap results as an additional finite-sample check.\footnote{The identifying cross-section contains 29 country clusters, so the finite-cluster $t$ and $F$ references use 28 denominator degrees of freedom. The restricted wild-cluster bootstrap follows \citet{mackinnon2017}.}

Our primary sample contains 29 African economies with both AERI and 2014 external-debt exposure. Figure~\ref{fig:external_identification} shows a delayed differential response. Four months after a one-standard-deviation tightening, a country one standard deviation more exposed before the sample records a 0.272-point larger cumulative AERI increase. The finite-cluster 90 percent interval is $[0.108,0.436]$ and $p=0.0086$. By month nine the differential reaches 0.384 points, with interval $[0.085,0.684]$ and $p=0.0373$. Restricted wild-cluster-bootstrap $p$-values are 0.042 and 0.034 at months four and nine, respectively. The response is delayed rather than contemporaneous, a pattern consistent with a gradual adjustment of the reported risk environment following tighter external financing conditions.

\exhibit{Figure_8.pdf}{AERI response to an externally identified U.S. monetary-policy tightening by predetermined external-debt exposure}{fig:external_identification}{The sample contains 29 African economies with AERI and 2014 external-debt exposure. The line reports $\hat\beta_h$ from equation~\eqref{eq:external_id}; shaded bands are pointwise 90 percent finite-cluster intervals using 28 denominator degrees of freedom. The shock is the standardised pure monetary-policy surprise from \citet{jarocinski2020}, and exposure is standardised $\log(1+\text{2014 external debt/GNI})$. The dependent variable is the cumulative change in AERI from month $t-1$ through horizon $h$. Country and year-month fixed effects and three lags of the outcome and interaction are included; standard errors are two-way clustered by country and month. The differential response is positive and statistically significant at months 4 and 9 under the finite-cluster reference.}

\begin{table}[htbp]\centering
\caption{Externally identified monetary-policy exposure validation}
\label{tab:external_identification}
\begin{singlespace}\scriptsize
\setlength{\tabcolsep}{3pt}
\begin{tabular}{@{}p{4.7cm}cccc@{}}
\toprule
Specification & $\hat\beta_4$ & $p_4$ & $\hat\beta_9$ & $p_9$\\
\midrule
Primary: log debt/GNI & 0.272 & 0.0086 & 0.384 & 0.0373\\
Raw debt/GNI & 0.203 & 0.0152 & 0.289 & 0.0556\\
95th-percentile-capped debt/GNI & 0.230 & 0.0284 & 0.347 & 0.0566\\
Exclude Mozambique & 0.265 & 0.0305 & 0.390 & 0.0610\\
African-country panel restricted to FOMC announcement months & 0.299 & 0.0078 & 0.354 & 0.0395\\
Alternative monetary-policy shock & 0.251 & 0.0068 & 0.420 & 0.0534\\
\midrule
Leave-one-country-out range & [0.214, 0.308] & 29/29 $p<0.10$ & [0.282, 0.455] & 29/29 $p<0.10$\\
\bottomrule
\end{tabular}
\end{singlespace}
\begin{singlespace}\footnotesize\textit{Notes:} The dependent variable is the cumulative change in AERI from month $t-1$ through horizon $h$. The shock is the standardized pure U.S. monetary-policy surprise from the Jaroci\'{n}ski--Karadi high-frequency decomposition, aggregated to the month. Predetermined exposure is 2014 external debt stocks as a share of GNI from the World Bank, transformed as $\log(1+x)$ and standardized in the primary specification. All regressions include country and year-month fixed effects, three lags of the AERI change and three lags of the shock--exposure interaction. Standard errors are two-way clustered by country and month. Reported $p$-values use a finite-cluster $t$ reference with degrees of freedom equal to the number of African-country clusters minus one. Restricted wild-cluster-bootstrap $p$-values for the primary month-four and month-nine coefficients are 0.042 and 0.034, respectively. In the leave-one-country-out exercise, 28 of 29 month-four estimates and 24 of 29 month-nine estimates also satisfy $p<0.05$.\end{singlespace}
\end{table}


Table~\ref{tab:external_identification} reports alternative exposure definitions, sample changes and a different monetary-policy shock. Every month-four and month-nine estimate remains positive and significant at the paper's 10 percent criterion under finite-cluster inference. The month-four coefficient also remains significant at 5 percent across all six reported variants. Month-nine estimates vary more in magnitude but retain a positive sign and 10 percent significance across all six variants. The result also remains when the time sample is restricted to FOMC announcement months.\footnote{This robustness specification changes the time observations only; the cross-section remains the same set of African economies.} In the leave-one-country-out calculations, the month-four coefficient ranges from 0.214 to 0.308, with all 29 estimates significant at 10 percent and 28 at 5 percent. At month nine it ranges from 0.282 to 0.455; all 29 are significant at 10 percent and 24 at 5 percent.

The falsification tests, reported in Supplementary Section~S13 and Table~S6, do not show the same pattern. Current monetary-policy surprises have little relationship with AERI changes dated one to six months earlier; finite-cluster $p$-values range from 0.156 to 0.807. The central-bank-information component of the same announcements produces no comparable debt-exposure pattern at months four and nine ($p=0.153$ and $p=0.125$). The joint future-shock interaction placebo yields $F(6,28)=1.884$ with $p=0.119$, and the restricted wild-cluster bootstrap gives $p=0.221$.\footnote{The corresponding large-sample Wald statistic is $\chi^2(6)=11.30$ with asymptotic $p=0.079$. With 29 country clusters, the finite-cluster and restricted wild-cluster-bootstrap references provide the more appropriate finite-sample benchmarks.} The 2014 debt ratio is predetermined but observational. It may be correlated with exchange-rate regimes, trade openness, institutional quality and other persistent country characteristics. The coefficient measures how the response to the external shock varies with pre-existing debt exposure; it does not by itself isolate a causal debt-transmission channel.

\FloatBarrier

\section{Conclusion}\label{sec:summary}
The paper constructs two news-based measures of economic risk for 54 African economies over 2015--2025. AERI aggregates graded article-level scores, the pillar indices retain the economic domain attached to those scores, and ARS provides a bounded summary built from the monthly risk signal. The more than 21 million retained article records show both common episodes and different country histories. The macroeconomic results are also directional rather than uniform: inflation-risk readings precede movements in headline inflation, exchange-rate depreciation is followed by higher AERI, stronger real activity is followed by lower AERI, and an externally identified U.S. monetary-policy tightening is followed by a larger AERI increase in African economies with greater predetermined external-debt exposure.

AERI, the pillar indices and ARS answer different questions. AERI records the amount of scored economic-risk content relative to the news flow. The pillar indices show which domains account for it. ARS indicates whether the monthly signal is narrow or spread across more of the risk dimensions used in the score. A sharp AERI increase with little movement in ARS therefore describes a different episode from one in which both measures rise.\footnote{Joint movements in AERI, ARS and the pillar indices distinguish concentrated episodes from broader cross-domain shifts and guide follow-up analysis.}

Cross-country comparison also has limits. Economic structure, external exposure, media depth, language and the availability of conventional high-frequency data differ across the sample. The common scoring rule fixes the definition of the measure, but country interpretation still depends on the source mix and the composition of the pillar scores.\footnote{The same headline value can reflect different pillar compositions, so country interpretation should use the pillar profile and coverage diagnostics.} Supplementary Section~S12 finds little within-country association between AERI and retained article volume. Newspaper archives can therefore add timely information where conventional data arrive slowly, but the index remains a measure of the observed news environment rather than a substitute for country-specific macroeconomic analysis.

\subsection{Limitations}\label{sec:limitations}
The measures are bounded by what appears in the retained news archive. AERI, the pillar indices and ARS describe reported economic risk in that information environment; they do not measure realised losses or crisis probabilities. Cross-pillar comparisons have the same limitation. A high financial-sector pillar mean, for example, means that financial-sector concerns are more prominent in the corpus, reflecting both underlying conditions and reporting practices.

Media coverage creates a separate source of uncertainty. Publication volume, editorial priorities, source continuity and language coverage differ across countries and over time. In the eligible sample, AERI has little within-country association with log article volume, so archive depth alone does not explain the index. That result does not establish source representativeness. The production dictionaries explicitly cover English, French and Portuguese, and reporting in other languages may be captured less completely. Exact-text deduplication removes identical reprints but leaves edited versions, translations and related stories as distinct records. The reported deduplication sensitivity applies to the descriptive series; the PVAR and external-shock specifications continue to use the original benchmark vintage.

The PVAR coefficients are conditional averages for the selected countries and depend on a researcher-imposed contemporaneous ordering. The main sign pattern appears under both the forward and feedback orderings, but the recursive system is not a structural causal design. We use it to describe dynamic relationships and rely on the external-shock specification for the separate identification exercise.

The monetary-policy shock is externally identified, but 2014 external-debt exposure is observational and may reflect persistent country differences. In the 29-country sample, we therefore read the result as an exposure-sensitive AERI response, not a causal debt-transmission estimate.
\printbibliography[title={References}]
\end{document}
